% Section 5.3, from lectures/L05/appendix/b_exercises.md.

\section{Exercises}
\label{c5:sec:exercises}\label{c5:app:b}

\begin{aside}
\textbf{How to work these.} The Reflection and Calculation exercises are answered on paper, from
the chapter. The Code exercises are programs for the board: build and flash them the way
\secref{c2:app:c} (Microchip Studio) or \secref{c2:app:d} (Linux) describes, or run them in the
simulator of Appendix~\ref{sim:ch} if you have no board yet.

\textbf{How to check your work.} The Reflection and Calculation exercises, and the questions some
Code exercises ask about their design, are answered in Appendix~\ref{app:answers}. The Code
exercises have no printed solution: the board tells you whether a program does what the exercise
describes. Try each exercise before you read its answer.
\end{aside}

% ------------------------------------------------------------------------------------------
\exerciseset{Concepts}
\label{c5:set:1}

\calcexercise{Tick at \textdiv256}
\label{c5:ex:1}
At 16\,MHz with a \textdiv256 prescaler, what's the resulting tick frequency and tick period? Show
your work.

\calcexercise{Overflows for 250\,ms}
\label{c5:ex:2}
Using your answer to Exercise~\ref{c5:ex:1}, how many overflow interrupts (8-bit timer) would be
needed to time a 250\,ms delay?

% ------------------------------------------------------------------------------------------
\exerciseset{Normal vs.\ CTC}
\label{c5:set:2}

\reflexercise{When to reach for CTC}
\label{c5:ex:3}
In your own words, explain when you'd reach for CTC mode instead of Normal mode with software
overflow-counting. Is there a period that's awkward to hit precisely in Normal mode but trivial in
CTC?

% ------------------------------------------------------------------------------------------
\exerciseset{Program}
\label{c5:set:3}

\codeexercise{Blinking at 2\,Hz}
\label{c5:ex:4}
Adapt the Timer 0 blink example of \secref{c5:a:4} to blink at 2\,Hz (250\,ms on, 250\,ms off)
instead of 1\,Hz. What value changes, and what does the new \code{OVF_MAX} work out to?

\reflexercise{When \code{volatile} matters}
\label{c5:ex:5}
The chapter's Timer 0 example declares \code{ovf_count} as \code{volatile}, and \secref{c5:a:4}
says that removing the qualifier there would change nothing. Explain why. Then
suppose \code{ovf_count} moved to file scope so that \code{main()} could wait for it, in a loop
such as \code{while (10U > ovf_count) {}}. What specific failure could removing \code{volatile}
now introduce, and why might it show up in one build and not in another?

\codeexercise{A 2-second period in CTC}
\label{c5:ex:6}
Configure Timer 1 in CTC mode for a period of exactly 2 seconds at 16\,MHz with a \textdiv1024
prescaler. What value goes in \code{OCR1A}? (Timer 1 is 16-bit, so check whether your answer
fits.)

\codeexercise{Three timers, one blink}
\label{c5:ex:7}
Capstone. Two LEDs on D9/D10 blink at a speed selected by whichever of three buttons (D11--D13)
was most recently pressed, each mapped to a different timer:
\begin{itemize}
  \item D11 selects Timer 0 (1000\,ms per toggle), D12 selects Timer 1 (500\,ms per toggle), D13
    selects Timer 2 (100\,ms per toggle).
  \item Pressing a button toggles \emph{that} timer's blinking on or off. If it was off, turn it on
    and turn the other two off (only one timer drives the LEDs at a time). If it was already the
    active timer, turn it off entirely and the LEDs go dark.
  \item All three buttons are on port B and share \code{PCINT0_vect}; detect which one was just
    pressed with the \code{prev}/\code{now} edge detection from \secref{c3:a:8}.
  \item Configure all three timers to tick at the same rate (\textdiv1024 prescaler,
    $\sim$0.064\,ms per tick) so the overflow-counting logic in \secref{c5:a:4} works identically
    for each. Timer 1 needs CTC mode with \code{OCR1A} set so its compare-match period matches
    Timer 0/2's overflow period, even though its 16-bit range could hit the target directly, so all
    three interrupt handlers end up the same shape.
  \item Define a small struct to hold \code{required_interrupts} and \code{executed_interrupts} for
    one timer (see \secref{c5:a:2} and \secref{c5:a:4} for the counting pattern), and use three
    instances, one per timer, instead of three sets of loose globals.
  \item Track the currently active timer in a single variable, using either an enum
    (\code{TIMER_NONE}, \code{TIMER_0}, \code{TIMER_1}, \code{TIMER_2}) or macros with timer
    indexes.
\end{itemize}
